\documentclass[10pt,a4paper]{article}
\usepackage[margin=2cm]{geometry}
\usepackage{fancyvrb,booktabs,enumitem,amsmath}
\fvset{fontsize=\scriptsize,frame=single}
\title{Lab Report -- Search and A* on the Warehouse Map}
\author{Apurva Verma}\date{}
\begin{document}\maketitle

\section*{Task 0: Search problem formulation}
\begin{center}\begin{tabular}{p{2.6cm}p{13cm}}\toprule
Component & Specification\\\midrule
State $S$ & a cell $(r,c)$ of the $9\times17$ grid that is not an obstacle\\
Actions $A$ & Up, Down, Left, Right\\
Transition $T$ & $T((r,c),a)=(r+\Delta r_a,c+\Delta c_a)$ if that cell is inside the grid and not \texttt{\#}\\
Initial state $s_0$ & position of \texttt{S}: $(1,1)$\\
Goal $G$ & $\{(7,15)\}$, the position of \texttt{G}\\
Cost $c$ & 1 per move, so path cost = number of moves\\\bottomrule\end{tabular}\end{center}
(a) A state needs only the robot's $(row,col)$; the map is fixed and need not be part of the state. (b) An action is invalid if it leaves the grid or enters a \texttt{\#} cell. (c) Yes, deterministic: each action in a state has exactly one outcome. (d) A solution is a sequence of valid actions leading from $s_0$ to the goal; an optimal solution has the fewest moves.

\section*{Task 1: Agent design}
\begin{enumerate}[leftmargin=*]
\item State: Python tuple \texttt{(row, col)}. \item Warehouse: list of strings converted to a list of lists; \texttt{S} and \texttt{G} located by scanning. \item Valid actions: for each of the four offsets, check bounds and \texttt{grid[r][c] != '\#'}. \item Goal test: \texttt{state == goal} when a state is taken from the frontier. \item Frontier (A*): min-priority queue of $(f,\text{tie-breaker},\text{state})$ with a dictionary $g$ of best-known costs; BFS uses a FIFO queue. \item Path: store a \texttt{parent} dictionary and walk back from the goal to \texttt{S}, then reverse.
\end{enumerate}
Reported on termination: found/not found, path, path length, states expanded (number of states removed from the frontier and expanded, excluding the goal itself).

\section*{Task 2: Prompts}
Full prompts are in the appendix and in \texttt{prompts.txt}. The first prompt followed the lab's example prompt with the exact map added and a request for a BFS version with the same interface. The code was produced with LLM assistance from the prompts in the appendix.

\section*{Task 3: Tests}
\begin{tabular}{p{3.6cm}p{12.5cm}}\toprule
Test & Result\\\midrule
1 Original map & found; length 40; 63 states expanded; path starts at S, ends at G, never enters \texttt{\#}, each step is one cell (asserted in code). Path: down the left column passage, along row 1 to col 5, down to row 5, east to col 13, up and back west along row 3, up to row 1 at col 7, east to col 15, down to G.\\
2 Trivial \texttt{\#SG\#\#} & found; length 1; 1 state expanded.\\
3 No solution & A* and BFS both report ``No solution found'' after expanding 9 states (the reachable region) and terminate, so no infinite loop thanks to the closed set.\\
4 Alternative paths (open $5\times5$ room, S and G at opposite corners) & many shortest paths exist; A* length 8 = BFS length 8 = Manhattan distance, so the returned path is shortest.\\\bottomrule\end{tabular}

\section*{Task 4: Where concepts appear in the code (\texttt{search\_agent.py})}
\begin{tabular}{p{3cm}p{13cm}}\toprule
Concept & Location\\\midrule
State & \texttt{(row, col)} tuples, everywhere (\texttt{parse} returns start/goal as tuples)\\
Action & \texttt{ACTIONS} dictionary (name $\to$ row/col offset)\\
Transition & \texttt{neighbours()}: applies the offset and returns the next state\\
Goal test & \texttt{if s == goal} after popping from the frontier in \texttt{astar}/\texttt{bfs}\\
$g(n)$ & dictionary \texttt{g}; \texttt{ng = g[s] + 1}\\
$h(n)$ & \texttt{manhattan(s, g)} (and variants \texttt{zero}, \texttt{euclidean}, \texttt{double\_manhattan})\\
$f(n)$ & \texttt{ng + h(n, goal)}, the first element of each heap entry\\
Frontier & \texttt{frontier} list used with \texttt{heapq} (A*); \texttt{deque} (BFS)\\
Visited states & \texttt{closed} set (A*); \texttt{parent} dictionary (BFS)\\
Path reconstruction & \texttt{rebuild(parent, s)}\\\bottomrule\end{tabular}

(a) A binary min-heap (\texttt{heapq}). (b) It pops the entry with smallest $f$; ties are broken by insertion order (a counter), or optionally by smaller $h$. (c) In \texttt{manhattan}, called when a neighbour is pushed. (d) Yes, explicitly: \texttt{ng + h(n, goal)}. (e) The \texttt{closed} set: a popped state already in it is skipped, and a neighbour is pushed only if its $g$ improves.

\section*{Task 5: BFS versus A* (same warehouse)}
\begin{center}\begin{tabular}{lcc}\toprule Measure & BFS & A*\\\midrule Solution found & yes & yes\\ Path length & 40 & 40\\ States expanded & 63 & 63\\\bottomrule\end{tabular}\end{center}
(a) Yes. (b) Yes, 40 each (identical path). (c) Neither: both expanded 63 states. (d) A* can expand fewer states because $h$ lets it ignore cells that look far from the goal. Here it does not help, because the warehouse is a winding maze with essentially one route: 63 of the 64 free cells (all but the goal) must be visited before the goal is reached by any algorithm, and Manhattan distance points through walls so it gives no useful guidance. On an open room the difference appears: with ties broken toward lower $h$, A* expanded 14 states versus 59 for BFS on a $6\times10$ room, and 8 versus 24 on the $5\times5$ room. Note that with the default insertion-order tie-breaking A* expanded the same 59 and 24 as BFS in open rooms, because every cell has the same $f$; tie-breaking is part of A*'s behaviour.

\section*{Task 6: Heuristic investigation}
Manhattan is appropriate because with four-directional unit-cost moves the true cost to the goal is at least the horizontal plus vertical distance (walls only add cost), so it never overestimates (admissible) and it is consistent.
\begin{center}\begin{tabular}{lccc|ccc}\toprule
 & \multicolumn{3}{c|}{Warehouse} & \multicolumn{3}{c}{Open $6\times10$ room}\\
Heuristic & found & length & expanded & found & length & expanded\\\midrule
Manhattan & yes & 40 & 63 & yes & 14 & 59\\
$h=0$ & yes & 40 & 63 & yes & 14 & 59\\
Euclidean & yes & 40 & 63 & yes & 14 & 59\\
$2\times$Manhattan & yes & 40 & 63 & yes & 14 & 14\\\bottomrule\end{tabular}\end{center}
(Default tie-breaking; the open room is an extra map, because the warehouse gives identical rows.) $h=0$ turns A* into uniform-cost search (same as BFS here). Euclidean is admissible but smaller than Manhattan, so it is less informed; in these runs it happened to expand the same number of states as Manhattan. Doubling Manhattan is inadmissible: on the open room it expanded only 14 states (greedy-like, fast) with the correct length, but on a small map with a dead-end corridor (\texttt{\#\#\#\#\#\#\#\#\#/\#S......\#/\#.\#.....\#/\#.\#...\#.\#/\#.\#\#..\#.\#/\#.....\#G\#}) the shortest path is 10, $2\times$Manhattan still finds 10, but $3\times$ and $10\times$ Manhattan return length 16 (and expand only 18 states versus 27). So a too-optimistic ($h=0$) heuristic costs time but keeps optimality, while an over-aggressive (inadmissible) one reduces work at the risk of a longer path.

\section*{Task 7: Evaluating the LLM-generated agent}
\textbf{Provenance:} the problem formulation and design (Tasks 0--1) follow the lab text; the data structures (\texttt{heapq}, \texttt{closed} set, tie-breaker counter), the code and the test harness were produced with LLM assistance. The A* and BFS implementations were accepted unchanged after the first run; an optional \texttt{prefer\_low\_h} tie-breaking flag and custom heuristics were then added for the experiments, and testing followed Tasks 3, 5 and 6.
\begin{enumerate}[leftmargin=*]
\item The state/transition/goal-test/path code and the closed-set logic were correct immediately (all four tests passed on the first run).
\item No correctness bug in the algorithm; one design issue: the heap needs a tie-breaker (otherwise tuples compare states) and the outcome depends on tie-breaking, which I had not anticipated.
\item By comparing A* against BFS in the open-room test, where both expanded identical numbers of states.
\item The tie-breaker counter (\texttt{itertools.count}) and ``stale entries'' in the heap needed explanation.
\item Yes: \texttt{prefer\_low\_h} and the heuristic arguments after the first version.
\item Test 4 (comparison with BFS as ground truth) and the no-solution test (termination).
\item No: the path ``looks plausible'' but only the BFS cross-check, obstacle check and unreachable test show it is shortest and safe.
\item A* is only as good as its heuristic; identical expansions to BFS on a maze show that a heuristic can be correct yet uninformative; tie-breaking matters; an inadmissible heuristic can return a longer path.
\end{enumerate}

\section*{Final Reflection}
\textbf{1.} Formulating the problem first fixes what a state, action, goal and cost are, which determines the data structures and what ``correct'' means; without it the code cannot be tested against anything. \textbf{2.} A* uses problem-specific knowledge, the heuristic $h(n)$, to order the frontier by $f=g+h$ rather than by depth or insertion order alone like blind search. \textbf{3.} An admissible heuristic keeps A* optimal; a larger admissible one prunes more; $h=0$ gives blind search; an inadmissible one (here $3\times$ Manhattan) can return longer paths; in the warehouse the heuristic gave no benefit because walls make Manhattan misleading. \textbf{4.} The LLM produced working, readable code and a test harness quickly and proposed the tie-breaker and closed-set structure; the design, experiments and interpretation were done in the same LLM-assisted session. \textbf{5.} Plausible output could hide non-optimal paths, wrong handling of unreachable goals (infinite loops), or expansion-count errors; here the tests and BFS ground truth are what establish correctness.

\section*{Appendix A: Prompts}
\VerbatimInput{prompts.txt}
\section*{Appendix B: Test output}
\VerbatimInput{results.txt}
\section*{Appendix C: Code}
\subsection*{search\_agent.py}\VerbatimInput{search_agent.py}
\subsection*{tests\_and\_experiments.py}\VerbatimInput{tests_and_experiments.py}
\end{document}
